\documentclass[10pt,twocolumn]{article}
\usepackage[T1]{fontenc}
\usepackage[utf8]{inputenc}
\usepackage{lmodern}
\usepackage[margin=1.8cm,top=2.4cm,bottom=2cm,columnsep=0.6cm]{geometry}
\usepackage{fancyhdr}
\usepackage{titlesec}
\usepackage{booktabs}
\usepackage{amsmath,amssymb,amsfonts}
\usepackage{microtype}
\usepackage{xcolor}
\usepackage{mdframed}
\usepackage{parskip}
\usepackage{enumitem}
\usepackage{hyperref}

\definecolor{rulered}{RGB}{180,30,30}
\definecolor{boxgray}{RGB}{245,245,245}
\definecolor{keywordgray}{RGB}{100,100,100}

\pagestyle{fancy}
\fancyhf{}
\fancyhead[L]{\textsc{\small Claude Mag}}
\fancyhead[C]{\small\textit{Machine Learning Research Digest}}
\fancyhead[R]{\small 2026-09-21}
\fancyfoot[C]{\small\thepage}
\renewcommand{\headrulewidth}{0.8pt}

\titleformat{\section}{\normalsize\bfseries\color{rulered}}{}{0pt}{}%
  [\vspace{-4pt}\color{rulered}\rule{\columnwidth}{0.4pt}]
\titlespacing{\section}{0pt}{8pt}{4pt}

\hypersetup{colorlinks=true,urlcolor=rulered,linkcolor=black}

\begin{document}
\setlength{\parindent}{0pt}
\setlength{\parskip}{4pt}

{\small\color{keywordgray}\textbf{Keywords:}
  multi-modal image generation \textbullet{}
  evaluation metrics \textbullet{}
  condition alignment \textbullet{}
  chain-of-evaluation}\par\vspace{4pt}

{\LARGE\bfseries\raggedright
  UFO}\par\vspace{4pt}

{\large\itshape\raggedright
  Atomized Chain-of-Evaluation Achieves 15\% Better Human-Correlation for
  Multi-Modal Image Generation by Evaluating All Conditions Simultaneously}\par\vspace{6pt}

{\small\textbf{By} Danning Zhang, Yijing Lin, Shuhan Zhuang, Mengqi Huang,
  Shaojin Wu, Shancheng Fang \& Zhendong Mao
  (University of Science and Technology of China)}\par\vspace{2pt}

{\small\color{keywordgray}arXiv:2609.12397 \textbullet{} ICML 2026}
\par\vspace{2pt}\noindent{\color{rulered}\rule{\columnwidth}{0.6pt}}\vspace{6pt}

Multi-modal image generation is evaluated by checking each condition
independently---a fundamentally wrong strategy when conditions interact
and conflict. \textit{UFO} introduces Atomized Chain-of-Evaluation:
decompose each condition into fine-grained Atomic Evaluation Units,
order them by dependency, call purpose-built verifiers, and aggregate.
On the new UFO-Bench (660 cases, 7 categories, 3 difficulty tiers),
UFO achieves Spearman~$\rho{=}0.689$ with human rankings---a 15.25\%
improvement over the best prior metric.

\begin{mdframed}[backgroundcolor=boxgray,linecolor=rulered,linewidth=1pt,
    innerleftmargin=6pt,innerrightmargin=6pt,
    innertopmargin=6pt,innerbottommargin=6pt]
\textbf{\small AT A GLANCE}\par\vspace{4pt}
\begin{description}[leftmargin=0pt,labelsep=4pt,itemsep=2pt,parsep=0pt,
    font=\small\bfseries]
  \item[Problem:] Existing multi-modal generation metrics evaluate each
    condition (text, reference image, layout, style, identity) in
    isolation, ignoring interactions that determine human preference.
  \item[Why it matters:] Subject-driven customisation and layout-
    conditioned generation are competitive research areas; unreliable
    metrics stall progress and misrank models.
  \item[Method:] Decompose each condition into Atomic Evaluation Units
    (AEUs); categorise by modality-relevance class; order into a chain;
    call structured (detector-based) or generic (MLLM) verifiers per AEU.
  \item[Result:] Spearman $\rho=0.689$ vs.\ $\rho=0.598$ (VIEScore);
    $+15.25\%$ average improvement. Largest gain on hard
    (conflicting-condition) tier.
  \item[Evidence:] UFO-Bench: 660 human-annotated cases, 5 raters each,
    Krippendorff $\alpha{=}0.73$; evaluated on CLIP-T, DINO-I, HPS-v2,
    VIEScore, and six other baselines.
\end{description}
\end{mdframed}

\section{The Big Picture}

Subject-driven text-to-image models accept a reference image (the
subject), a text prompt, a layout specification, and optionally a style
or identity condition---all simultaneously. The output is evaluated for
how well it aligns with all provided conditions at once.

Current metrics fail this task. CLIP-T measures text-image cosine
similarity; DINO-I measures instance identity similarity; they are then
averaged into a scalar. But this average hides condition priority: if a
prompt says ``make the dog hold a red umbrella'' and the model produces
a dog without an umbrella, no averaging of text and identity scores
captures this failure correctly. Humans, presented with the same output,
immediately rank it below a generation that at least includes the umbrella.

UFO treats evaluation as a structured reasoning problem: enumerate what
must be verified, verify each part with the best available tool, and
combine respecting priority and dependency.

\section{Why Existing Approaches Fall Short}

\textbf{Embedding-based metrics} (CLIP-T, DINO-I, CLIP-I) measure
condition-specific cosine similarity in a shared embedding space.
Their scores are invariant to fine-grained compositional failures;
a tiger wearing a hat and a tiger without a hat may have nearly
identical CLIP scores with the same text prompt.

\textbf{MLLM-based metrics} (VIEScore, Pick-Score) use a multimodal
language model as a judge but ask it to evaluate all conditions
holistically in a single prompt. The model must internally reason
about interactions it was not specifically trained to handle; its
judgments are inconsistent on conflicting-condition cases.

\textbf{Human Preference Score (HPS-v2)} is trained on general human
preferences that reflect aesthetics and realism as much as condition
alignment, making it poorly calibrated for the specific task of
measuring omni-condition satisfaction.

\section{Core Mechanism}

\textbf{AEU decomposition.} For a set of conditions $C = \{c_1, \ldots, c_K\}$
(text prompt $c_T$, reference image $c_I$, layout $c_L$, style $c_S$, \ldots),
each condition is parsed into binary claims about one attribute:
\[
\{u_{k,j}\} = \text{Parse}(c_k), \quad
u_{k,j} : \text{``attribute } a_{k,j} \text{ of } c_k \text{ satisfied?''}
\]

\textbf{Chain ordering.} AEUs are sorted topologically over a dependency
graph $G_C$ where edges encode precedence (subject identity must be
located before subject pose can be verified). Mandatory conditions come
first; preferred and contextual follow.

\textbf{Specialised verifiers.}
\begin{itemize}[itemsep=1pt,parsep=0pt]
  \item \emph{Text attributes:} MLLM with targeted yes/no prompts.
  \item \emph{Identity conditions:} ArcFace (persons) or DINOv2
    cosine similarity (objects).
  \item \emph{Geometric conditions:} DWPose (body pose), Grounding-DINO
    + SAM2 (object placement), DepthAnything (depth layout).
  \item \emph{Style conditions:} CLIP style embedding vs.\ reference.
\end{itemize}

\textbf{Weighted aggregation.} Mandatory AEUs have weight 1.0,
preferred 0.7, contextual 0.3:
\[
\text{OCA}_{\text{UFO}}(x,C) = \sum_{k,j} r_{k,j}\cdot V_{k,j}(x, c_k, a_{k,j}).
\]

\section{Technical Formulation}

Let $\mathcal{U} = [u_1, \ldots, u_n]$ be the ordered chain of AEUs
and $V_i$ the verifier for AEU $u_i$. The OCA score is:
\[
\text{OCA}(x,C) = \frac{\sum_{i=1}^n r_i V_i(x)}{\sum_{i=1}^n r_i},
\quad r_i \in \{1.0, 0.7, 0.3\}.
\]

Spearman correlation with human rankings:
\[
\rho = 1 - \frac{6\sum_i d_i^2}{n(n^2-1)},
\]
where $d_i$ is the rank difference between UFO scores and human
preference rankings over the $n=660$ UFO-Bench cases.

For text prompt $c_T$ with $P$ phrases, each phrase $p$ is parsed into
an AEU with attribute type (object, property, spatial, action, count)
and relevance class. The dependency graph adds edges from objects to
their attributes (the object must be localised before its colour or
pose can be verified).

\section{Learning or Inference Procedure}

UFO is training-free. Inference steps:

\begin{enumerate}[itemsep=2pt,parsep=0pt]
  \item Parse all conditions $C$ into AEUs; extract attribute types and
    relevance classes.
  \item Build dependency graph $G_C$; topologically sort into chain $\mathcal{U}$.
  \item For each AEU $u_i$: call verifier $V_i$ with the generated image
    $x$ and condition $c_k$.
  \item Compute weighted aggregation $\text{OCA}_{\text{UFO}}(x,C)$.
\end{enumerate}

New condition types are added by registering an AEU parser and
verifier function---no retraining. Total evaluation time per image is
approximately 4~seconds on a single GPU.

\section{What the Guarantee Says}

UFO provides a theoretical lemma: under conditional independence between
conditions, the chain-of-evaluation score is an unbiased estimator of
the true omni-condition alignment. For the general (dependent) case,
the mandatory-first ordering minimises the expected rank correlation
error by prioritising the conditions humans weight most heavily, which
the paper verifies empirically across annotation tiers.

\section{Experimental Findings}

Rank correlation with human preferences on UFO-Bench (660 test cases,
5 annotators each):

\begin{center}
\small
\begin{tabular}{lccc}
\toprule
Metric & Spearman $\rho$ & Pearson $r$ & Kendall $\tau$ \\
\midrule
CLIP-T & 0.421 & 0.398 & 0.312 \\
DINO-I & 0.389 & 0.374 & 0.288 \\
HPS-v2 & 0.521 & 0.508 & 0.401 \\
VIEScore & 0.598 & 0.581 & 0.463 \\
\textbf{UFO} & \textbf{0.689} & \textbf{0.672} & \textbf{0.537} \\
\bottomrule
\end{tabular}
\end{center}

By difficulty tier: Easy $\rho=0.741$; Medium $\rho=0.693$;
Hard $\rho=0.701$. UFO's advantage is largest on hard cases, where
conflicting conditions are the norm and independent metrics collapse.

\section{Ablations and Interpretation}

\textbf{MLLM for all AEUs (no structured verifiers):} $\rho$ drops to
0.631; geometric and identity AEUs are specifically hurt, confirming
that dedicated detectors are needed for these types.

\textbf{Random AEU ordering:} $\rho=0.657$; mandatory-first ordering
resolves conflicting conditions more consistently.

\textbf{Uniform weights ($r_i{=}1$ for all):} $\rho=0.672$; graded
weights align with human priority on medium and hard tiers.

\textbf{Sentence-level vs.\ phrase-level AEUs:} Phrase-level AEUs
improve $\rho$ by 0.042; sub-phrase (word-level) adds only 0.003
at 2$\times$ compute.

\textbf{Model comparison on UFO-Bench:} FLUX.1 outranks SDXL by
0.08 UFO score on hard cases but only 0.02 by CLIP-T average,
illustrating UFO's superior discriminative power for state-of-the-art
models that are close to each other on simple metrics.

\section*{Reference}

Danning Zhang, Yijing Lin, Shuhan Zhuang, Mengqi Huang, Shaojin Wu,
Shancheng Fang, Zhendong Mao.
\textbf{UFO: Chain-of-Evaluation for Omni-Condition Alignment in
Multi-Modal Image Generation.}
arXiv:2609.12397, September 2026. ICML 2026.
\url{https://arxiv.org/abs/2609.12397}

\end{document}
